\section{Background and Related Work}
\label{sec:bg}

\subsection{AIMer v3 and AIM3}

AIMer proves knowledge of a secret input whose image under a public one-way
function is contained in the public key.  The proof is obtained by simulating
an MPC protocol among $N$ virtual parties and applying the Fiat--Shamir
transform.  This follows the MPCitH paradigm of Ishai et
al.~\cite{ishai2007zero} and the large-field checking techniques used in
BN++~\cite{kales2022efficient}.

Version 3 changes the underlying one-way function from AIM2 to AIM3.  For a
field element $x\in\gftwo{n}$, AIM3 uses the enhanced Mersenne S-box
\begin{equation}
  S[e](x)=x^{2^e-1}+x^{2^n-2}.
\end{equation}
An IV-derived invertible affine map is applied before the input S-boxes, and
their outputs pass through an output affine layer.  The implicit S-box relation
is valid only for nonzero inputs, so key generation rejects a candidate when an
S-box input is zero~\cite{cho2026aimer3}.  The signature layer also uses
AIM3-specific linear generation, MPC simulation, and S-box-output algorithms.

Table~\ref{tab:params} lists the six parameter sets used in this work.  The
\emph{f} sets use fewer parties and more repetitions, whereas the \emph{s} sets
use 256 parties and fewer repetitions to obtain smaller signatures.

\begin{table}[t]
\centering
\caption{AIMer v3 parameter sets.  $n$ is the field and security size, $N$ is
the number of MPC parties, and $\tau$ is the number of repetitions.}
\label{tab:params}
\small
\setlength{\tabcolsep}{5.5pt}
\begin{tabular}{lrrrrrr}
\toprule
Set & $n$ & Input S-boxes & $N$ & $\tau$ & sk (B) & Signature (B)\\
\midrule
128f & 128 & 2 & 16  & 33 & 48 & 6,944\\
128s & 128 & 2 & 256 & 17 & 48 & 4,704\\
192f & 192 & 2 & 16  & 49 & 72 & 15,408\\
192s & 192 & 2 & 256 & 25 & 72 & 10,320\\
256f & 256 & 3 & 16  & 65 & 96 & 31,360\\
256s & 256 & 3 & 256 & 33 & 96 & 20,224\\
\bottomrule
\end{tabular}
\end{table}

The cost relevant to this work appears in two repeated structures.  First,
party commitments and random tapes are derived from SHAKE128 at security level
1 and SHAKE256 at levels 3 and 5.  Second, the affine layers, Frobenius powers,
and multiplication checks repeatedly operate on independent party shares over
$\gftwo{n}$.  Parallelizing parties changes neither the field equations nor the
proof transcript order.

\subsection{AIM, AIM2, and AIMer Implementations}

The original AIM primitive was designed to minimize multiplicative complexity
inside an MPCitH proof~\cite{kim2023aim}.  Lee et al. optimized its Mersenne
operations and linear layer and reported a substantial improvement over the
reference code~\cite{lee2023high}.  Algebraic analyses then motivated AIM2,
which modified the primitive and its parameters~\cite{kim2023aim2}.  More
recent attacks on full AIM2 motivated the transition to AIM3
~\cite{biryukov2026magic,sun2025high}.

Prior AIMer work also improved the proof-system layer.  Kim et al. applied
relaxed vector commitments and GGM-tree techniques to AIMer v2, reducing both
signature size and execution time~\cite{kim2025relaxed}.  These changes are
orthogonal to the low-level SIMD kernels considered here.

The official implementations of earlier AIMer versions included AVX2 code,
and our previous AIM2 work added AVX-512 field and Keccak
paths~\cite{lee2026aimeravx512}.  The AIM3 specification, released in August
2026, defines the new algorithms but does not report an AVX2/AVX-512
comparison or detailed performance tables~\cite{cho2026aimer3}.  Consequently,
it remained unclear which earlier optimization principles survive the
algorithm change.

\subsection{SIMD Building Blocks}

Keccak naturally supports parallel-instance SIMD because independent states
execute the same round function~\cite{bertoni2013keccak}.  For binary-field
arithmetic, PCLMULQDQ computes a carry-less $64\times64$ product, while
VPCLMULQDQ applies the operation independently to the four 128-bit lanes of a
ZMM register~\cite{drucker2018fast}.  AVX-512 additionally provides 32 vector
registers, ternary Boolean logic, lane rotations, and masked memory operations.
AVX-512VL exposes many of these operations on XMM and YMM operands
~\cite{anderson2018enhanced}.  Our implementation combines these known
building blocks with the party structure of AIMer v3.
